% Optional compact method text for Task 5; requires amsmath.
% Describes the exact CMV-positive-supporting selection used in the new plots.
We examine subsets selected in the CMV-positive direction. For CellCNN, a cell is selected
if any filter with positive output contrast exceeds half its original
training-reference maximum. For SVM, it must belong to the full donor's top
$\lceil 0.01N_d\rceil$ cell margins and exceed the saved inner-validation
threshold. For Citrus, it must map to a training cluster with a positive
classifier coefficient. The corresponding scores and selection rules are
\[
\begin{aligned}
q^{\mathrm{CNN}}_{is}
  &= \max_{f:\,\Delta_{fs}>0,\ M_{fs}>0} r_{ifs}/M_{fs},
  &q^{\mathrm{CNN}}_{is}&>0.5,\\
q^{\mathrm{SVM}}_{is}
  &= m_{is}-\tau_s\quad(i\in I_{ds}),
  &q^{\mathrm{SVM}}_{is}&>0,\\
q^{\mathrm{Citrus}}_{is}
  &= \max\bigl(\{0\}\cup\{\beta_{js}:i\in A_{js}\}\bigr),
  &q^{\mathrm{Citrus}}_{is}&>10^{-10}.
\end{aligned}
\]
Here $M_{fs}$ is the training-reference maximum and $\Delta_{fs}$ the
CMV-positive minus CMV-negative output weight; the CellCNN maximum is zero
if no filter is eligible. $I_{ds}$ contains the donor's top-scoring cells;
SVM cells outside it are unselected. $A_{js}$ is Citrus membership inherited
from the nearest original training cell in ArcSinh marker space.
Each cell is counted at most once per split, including overlapping filters
or clusters. The plotted frequency divides its positive selections by all
outer test appearances of its donor; empty selections remain in the
denominator. These rules describe model-associated subsets, not causal
cell-removal effects.
